% ─────────────────────────────────────────────────────────────────────────────
\section{Tools, Platforms, and Experiment Tracking}
\label{sec:ch5_tools_platforms}
% ─────────────────────────────────────────────────────────────────────────────

The experiments in this report were implemented in Python using a PyTorch-based
training pipeline. The codebase is organized around a configurable experiment
script that defines dataset construction, model initialization, adversarial attack
generation, training objectives, checkpoint selection, and fixed
evaluation scripts in one place. This organization is important because the compared methods differ
mainly in their grouping and weighting strategies, while the dataset split,
architecture, source attacks, evaluation attacks, optimizer, and checkpointing
rules are kept fixed.

Because reproducibility depends on both the training code and the analysis stack,
Table~\ref{tab:ch5_tools_platforms} records the tools and platforms that form the
experimental workflow. Google Colab and Google Drive are used for
interactive execution and persistent storage of checkpoints, cached adversarial
examples, figures, and result files. PyTorch and torchvision provide the core
deep-learning infrastructure, including model definitions, data loading, and
optimization. The attack implementations are built primarily with TorchAttacks,
with AutoAttack-$\ell_\infty$ and DDN-$\ell_2$ handled through their dedicated
packages when required. scikit-learn is used for K-means clustering in
AttackDRO++, while NumPy and pandas are used to aggregate seed-level and
attack-level results.

\begin{table}[H]
\centering
\caption{Tools and platforms used in the experimental workflow. The table lists only tools that are part of the implemented training, evaluation, logging, or reporting pipeline.}
\label{tab:ch5_tools_platforms}
\footnotesize
\setlength{\tabcolsep}{4pt}
\renewcommand{\arraystretch}{1.08}
\begin{tabularx}{\textwidth}{@{}
>{\raggedright\arraybackslash}p{0.24\textwidth}
>{\raggedright\arraybackslash}p{0.28\textwidth}
>{\raggedright\arraybackslash}X
@{}}
\toprule
Tool / platform & Role & Use in this report \\
\midrule

Google Colab
& Interactive execution environment
& Used to run training and evaluation notebooks with GPU support. \\

Google Drive
& Persistent experiment storage
& Stores checkpoints, cached adversarial examples, result tables, figures, and intermediate analysis files. \\

Python
& Main implementation language
& Used for the full training, evaluation, analysis, and plotting pipeline. \\

PyTorch
& Deep-learning framework
& Implements model training, automatic differentiation, mixed-precision execution, and checkpoint loading. \\

torchvision
& Dataset and model utilities
& Provides CIFAR-10 loading, image transforms, and ResNet backbones adapted for small images. \\

TorchAttacks
& Adversarial attack library
& Provides most gradient-based and optimization-based attacks used in training and evaluation. \\

AutoAttack package
& Standardized robustness stress test
& Used for AutoAttack-$\ell_\infty$ evaluation with perturbation budget $8/255$. \\

adv-lib
& DDN attack implementation
& Used for DDN-$\ell_2$ adversarial example generation. \\

scikit-learn
& Clustering utilities
& Provides MiniBatch K-means for cluster discovery in AttackDRO++. \\

NumPy and pandas
& Result processing
& Used to aggregate per-seed, per-attack, per-class, and graybox transfer results. \\

matplotlib and SciencePlots
& Visualization
& Used to generate plots and publication-style figures for analysis. \\

Weights \& Biases
& Experiment tracking
& Logs configurations, training curves, validation metrics, group weights, and diagnostic figures when enabled. \\

\bottomrule
\end{tabularx}
\end{table}

Experiment tracking is treated as part of the reproducibility protocol rather than
as a separate result source. Each run records the method name, seed, dataset,
backbone, source attacks, evaluation attacks, and key hyperparameters. During
training, the pipeline logs loss, adversarial accuracy, validation robustness,
learning rate, checkpoint state, and group-weight diagnostics when applicable.
Representative Weights \& Biases (W\&B) tracking exports are provided in
Appendix~\ref{app:logs}.

All test-set claims in Chapter~6 are based on saved checkpoints that are loaded
after training and evaluated using fixed evaluation scripts. The main result
tables are produced from structured result files rather than from manually copied
console logs. This separation reduces the risk of mixing validation, test, whitebox, graybox, and diagnostic metrics. It
also makes it possible to reuse the same result files for aggregate metrics,
per-attack analysis, per-class analysis, and paired statistical comparisons.

\paragraph{Summary.}
The setup summary in Table~\ref{tab:ch5_setup_summary} serves as the reproducibility
contract for the result claims in Chapter~6 by collecting the dataset, model, attack
suite, repeated-seed protocol, and analysis stack in one place.

\begin{table}[H]
\centering
\caption{Summary of the experimental setup used in this report. Unless otherwise stated, evaluations use CIFAR-10, ResNet-18, and the eight non-AutoAttack attacks for Mean(8), with AutoAttack-$\ell_\infty$ reported separately.}
\label{tab:ch5_setup_summary}
\footnotesize
\setlength{\tabcolsep}{4pt}
\renewcommand{\arraystretch}{1.08}
\begin{tabularx}{\textwidth}{@{}
>{\raggedright\arraybackslash}p{0.28\textwidth}
>{\raggedright\arraybackslash}X
@{}}
\toprule
Component & Configuration \\
\midrule

Dataset
& CIFAR-10 with 40{,}000 training images, 10{,}000 validation images, and 10{,}000 test images. \\

Backbone
& ResNet-18 adapted for CIFAR-10 resolution. \\

Main Chapter~6 compared methods
& PGD-AT, DDN-AT, Multi-AT, and AttackDRO++; AttackDRO and Cluster-DRO are intermediate design steps described in Chapter~\ref{chap:methodology}. \\

Source attacks
& PGD-$\ell_\infty$ and DDN-$\ell_2$. \\

Evaluation attacks for Mean(8)
& FGSM-RS, PGD-$\ell_\infty$, TPGD, MI-FGSM, PGD-$\ell_2$, DDN-$\ell_2$, DeepFool-$\ell_2$, and CW-$\ell_2$. \\

Separate stress test
& AutoAttack-$\ell_\infty$, reported separately from Mean(8). \\

Primary aggregate metric
& Mean(8): mean robust accuracy over the eight non-AutoAttack evaluation attacks. \\

Secondary metrics
& Worst(8), seen-source robustness, held-out $\ell_\infty$ robustness, held-out $\ell_2$ robustness, per-attack accuracy, per-class accuracy, and graybox transfer accuracy. \\

Repeated-seed protocol
& Aggregate whitebox tables use paired random-seed comparisons; class-wise and graybox panels use five random seeds per method. \\

Graybox protocol
& Four method families are evaluated across five random seeds each, giving 20 random-seed checkpoints and 400 surrogate--target pairs. \\

Statistical analysis
& Paired random-seed comparisons with Wilcoxon signed-rank tests, paired $t$-tests, bootstrap confidence intervals, effect sizes, and Holm--Bonferroni correction. \\

Implementation stack
& PyTorch, torchvision, TorchAttacks, AutoAttack package, adv-lib, scikit-learn, NumPy, pandas, matplotlib, SciencePlots, W\&B, Google Colab, and Google Drive. \\

\bottomrule
\end{tabularx}
\end{table}

This chapter has defined the experimental setting used throughout the remainder of
the report: the dataset, model, attacks, training configurations, hyperparameters,
statistical protocol, evaluation metrics, graybox transfer design, and
implementation tools. These definitions provide the reproducibility contract for
the empirical claims that follow. Chapter~\ref{chap:results_analysis} now presents the results obtained under this protocol and analyzes how the proposed method behaves under whitebox, graybox, per-class, and ablation settings.
